\documentclass[a4paper,14pt]{extarticle}
\usepackage{fontspec}
\IfFontExistsTF{Times New Roman}{\setmainfont{Times New Roman}}{\setmainfont{Liberation Serif}}
\IfFontExistsTF{Courier New}{\setmonofont{Courier New}}{\setmonofont{Liberation Mono}}
\usepackage[russian]{babel}
\usepackage[a4paper,left=30mm,right=15mm,top=20mm,bottom=20mm,footskip=10mm]{geometry}
\usepackage{graphicx}
\usepackage{longtable}
\usepackage{array}
\usepackage{fancyhdr}
\usepackage{tocloft}
\usepackage{fvextra}
\usepackage[hidelinks]{hyperref}

\fontsize{14}{18}\selectfont
\setlength{\parindent}{1.25cm}
\setlength{\parskip}{0pt}
\setlength{\baselineskip}{18pt}
\raggedbottom
\sloppy

\pagestyle{fancy}
\fancyhf{}
\fancyfoot[R]{\thepage}
\renewcommand{\headrulewidth}{0pt}

\renewcommand{\cftsecfont}{\normalfont}
\renewcommand{\cftsubsecfont}{\normalfont}
\renewcommand{\cftsecpagefont}{\normalfont}
\renewcommand{\cftsecleader}{\cftdotfill{\cftdotsep}}
\renewcommand{\cftsubsecleader}{\cftdotfill{\cftdotsep}}
\setlength{\cftbeforesecskip}{0pt}
\setlength{\cftbeforesubsecskip}{0pt}

\newcommand{\reportchapter}[2]{\clearpage\noindent\hspace{\parindent}\textbf{#1 #2}\par\vspace{\baselineskip}\addcontentsline{toc}{section}{#1 #2}}
\newcommand{\reportsection}[2]{\noindent\hspace{\parindent}\textbf{#1 #2}\par\vspace{\baselineskip}\addcontentsline{toc}{subsection}{#1 #2}}
\newcommand{\reportsubsection}[2]{\noindent\hspace{\parindent}\textbf{#1 #2}\par\vspace{\baselineskip}}
\newcommand{\reportplain}[1]{\clearpage\begin{center}\textbf{#1}\end{center}\vspace{\baselineskip}\addcontentsline{toc}{section}{#1}}
\newcommand{\reportappendix}[2]{\clearpage\begin{center}\textbf{ПРИЛОЖЕНИЕ #1}\\\textbf{#2}\end{center}\vspace{\baselineskip}\addcontentsline{toc}{section}{ПРИЛОЖЕНИЕ #1 #2}}

\begin{document}


\begin{titlepage}
\centering
Министерство образования Республики Беларусь\par
\vspace{\baselineskip}
Учреждение образования\par
БЕЛОРУССКИЙ ГОСУДАРСТВЕННЫЙ УНИВЕРСИТЕТ\par
ИНФОРМАТИКИ И РАДИОЭЛЕКТРОНИКИ\par
\vspace{\baselineskip}
\vspace{\baselineskip}
\noindent Факультет\hspace{1em}<факультет>\par
\vspace{\baselineskip}
\noindent Кафедра\hspace{1em}<кафедра>\par
\vspace{\baselineskip}
\noindent Дисциплина:\hspace{1em}Тестирование программного обеспечения\par
\vspace{\baselineskip}
\vspace{\baselineskip}
ОТЧЁТ\par
по лабораторной работе\par
на тему\par
\vspace{\baselineskip}
\textbf{РАЗРАБОТКА ТЕСТ-КЕЙСОВ}\par
\vspace{\baselineskip}
\vspace{\baselineskip}
\vspace{\baselineskip}
\vspace{\baselineskip}
\noindent Студент:\hfill <И.О. Фамилия студента>\par
\vspace{\baselineskip}
\vspace{\baselineskip}
\noindent Руководитель:\hfill <И.О. Фамилия руководителя>\par
\vspace{\baselineskip}
\vspace{\baselineskip}
\vspace{\baselineskip}
\vspace{\baselineskip}
\vspace{\baselineskip}
\vspace{\baselineskip}
\vspace{\baselineskip}
\vspace{\baselineskip}
Минск 2026\par
\end{titlepage}

\reportchapter{1}{ЦЕЛЬ РАБОТЫ}

Изучить структуру тест-кейса, освоить эквивалентное разбиение и анализ граничных значений, оформить и выполнить тест-кейсы для формы заявки на раннее подключение веб-приложения DSLKit.


\reportchapter{2}{ТЕОРЕТИЧЕСКИЕ СВЕДЕНИЯ}

Теория дана по книге С. С. Куликова «Тестирование программного обеспечения. Базовый курс» (EPAM Systems, 2015–2025), главы 2.4 и 2.7.2.


Тест-кейс — набор входных данных, условий выполнения и ожидаемых результатов, разработанный для проверки определённого свойства или поведения программы. Если не указаны входные данные, условия или ожидаемый результат, либо неясна цель, тест-кейс считается плохим. Тест-кейсы обеспечивают повторяемость проверок, передачу знаний в команде, регрессионное тестирование и улучшение требований.


Класс эквивалентности — набор данных, обрабатываемых одинаково и дающих одинаковый результат. Граничное условие — значение на границе классов эквивалентности. Для диапазона проверяют последнее допустимое и первое недопустимое значения, а также «опасные» значения вроде нуля и пустой строки. Позитивные тест-кейсы берут из допустимых классов, негативные из недопустимых.


\reportchapter{3}{ХОД РАБОТЫ}

\reportsection{3.1}{Выбор use case и требования}

Тестируется вариант использования UC-01, описанный ниже. Проверку полей выполняет функция validateWaitlist (файл site/js/utils/validation.js), которая возвращает список неверных полей. Проект развёрнут локально командой npm start.


*\textit{Название.}*


– UC-01. Подать заявку в лист ожидания (форма «Записаться в бету», главная страница DSLKit).



*\textit{Идентификатор, версия, приоритет.}*


– UC-01, версия 1.0, приоритет High: единственный способ оставить контакты для беты.



*\textit{Краткое описание цели и задачи.}*


– Посетитель оставляет контактные данные и сведения о проекте, чтобы команда DSLKit могла прислать ему приглашение в бету после запуска. Задача формы: принять только корректную заявку и сохранить её.



*\textit{Связанные требования.}*


– R4 «Форма заявки и валидация» (лаб. 3, раздел 1). Реализация: site/index.html (форма \#waitlist-form), site/js/components/waitlist-form.js, site/js/utils/validation.js, site/thanks.html и site/en/thanks.html.



*\textit{Подзадачи.}*


– UC-01.1 Показ формы и подсказок к полям (шаги 1–6);

– UC-01.2 Очистка и проверка значений по П1–П6 (шаги 8–9);

– UC-01.3 Показ ошибок, фокус и повторная проверка (A1, A2);

– UC-01.4 Сохранение заявки в LocalStorage (шаг 10, A4, E1);

– UC-01.5 Переход на страницу благодарности (шаг 11).



*\textit{Связанные задачи и use case.}*


– UC-00 «Прочитать политику конфиденциальности»: ссылка в тексте согласия ведёт на privacy.html (шаг 6);

– Лаб. 2: тест-кейсы TC-WL-001…022;

– Лаб. 3: проверки SM-07, CP-35, CP-36;

– Дефекты реализации: D-1 (повторные заявки), D-2 (потеря заявки при сбое хранилища).



*\textit{Соответствие STAR.}*


– Situation (ситуация): предусловия 1–2;

– Task (задача): краткое описание цели и задачи;

– Action (действия): основной поток, шаги 1–11, и альтернативные потоки A1–A4;

– Result (результат): постусловия и исключительные ситуации E1–E2.



*\textit{Основной актор.}*


– Посетитель сайта (неавторизованный пользователь). Внешних систем в сценарии нет.



*\textit{Предусловия.}*


– 1. Открыта главная страница сайта (русская или английская версия);

– 2. Форма «Записаться в бету» отображается, все поля пусты, поле «Размер команды» содержит 0.



*\textit{Основной поток событий.}*


– 1. Посетитель вводит имя в поле «Имя»;

– 2. Посетитель вводит адрес в поле «Электронная почта»;

– 3. Посетитель выбирает значение в списке «Нотация»;

– 4. Посетитель вводит целое число в поле «Размер команды»;

– 5. Посетитель при желании вводит текст в поле «Комментарий»;

– 6. Посетитель отмечает флажок согласия;

– 7. Посетитель нажимает кнопку «Отправить»;

– 8. Система удаляет пробелы по краям у имени, почты и комментария;

– 9. Система проверяет все поля по правилам П1–П6 (ниже);

– 10. Система сохраняет заявку в LocalStorage браузера под ключом dslkit.waitlist: имя, почта, нотация, размер команды, комментарий, время создания;

– 11. Система открывает страницу благодарности thanks.html.



*\textit{Правила проверки полей.}*


– П1. Имя: 2–60 символов, начинается с буквы, далее буквы, пробел, дефис, апостроф;

– П2. Почта: вид name@domain.zone, домен верхнего уровня не короче двух символов, длина не более 120;

– П3. Нотация: одно из значений idef0, dfd, erd, bpmn, uml, petri;

– П4. Размер команды: целое число от 0 до 500 включительно;

– П5. Комментарий: не более 500 символов, поле необязательное;

– П6. Согласие: флажок обязательно отмечен.



*\textit{Альтернативные потоки.}*


– A1. Одно или несколько полей не прошли проверку (шаг 9): под каждым неверным полем появляется сообщение об ошибке, фокус переходит на первое неверное поле, заявка не сохраняется. После исправления поля и ухода из него сообщение исчезает. Сценарий продолжается с шага 7;

– A2. Согласие не отмечено (шаг 9): применяется A1 для поля согласия;

– A3. Посетитель закрывает страницу или уходит со страницы до шага 7: данные не сохраняются, состояние системы не меняется;

– A4. Повторная заявка с тем же адресом почты (шаг 10): заявка принимается как новая отдельная запись, сообщения о дубликате нет.



*\textit{Постусловия.}*


– Успех: в LocalStorage на одну запись больше, открыта страница thanks.html;

– Ошибка проверки (A1, A2): записей не добавилось, посетитель остаётся на главной странице, введённые значения сохранены в полях.



*\textit{Исключительные ситуации.}*


– E1. Браузер не позволяет записать в LocalStorage (хранилище переполнено или запрещено, шаг 10): система всё равно открывает страницу благодарности, заявка не сохраняется, сообщения об ошибке нет. Это известное ограничение текущей версии (дефект D-2 лабораторной работы №1);

– E2. Одно из полей имеет недопустимый тип данных (например, текст в поле «Размер команды»): обрабатывается как A1 для этого поля.



Трассировка: правило П1 проверяют TC-WL-002…007, П2 008…011, П3 018, П4 012…017, П5 019 и 020, П6 021, поток A1 022.


Требования к полям приведены в таблице ниже.


\noindent Таблица 3.1 – Требования к полям формы\par\nopagebreak
\begin{longtable}{|>{\raggedright\arraybackslash}p{0.490\textwidth}|>{\raggedright\arraybackslash}p{0.490\textwidth}|}
\hline
Поле & Правило \\ \hline
\endhead
name & Начинается с буквы; буквы, пробел, апостроф, дефис; длина от 2 до 60 \\ \hline
email & Формат адреса с доменом верхнего уровня; длина не более 120 \\ \hline
notation & Одно из: idef0, dfd, erd, bpmn, uml, petri \\ \hline
team & Целое число от 0 до 500 (допустима строка из цифр) \\ \hline
comment & Не более 500 символов, необязательное \\ \hline
consent & Только true \\ \hline
\end{longtable}
\vspace{\baselineskip}

\reportsection{3.2}{Классы эквивалентности и границы}

Для каждого поля выделены допустимый и недопустимые классы; границы проверяются парами «последнее допустимое и первое недопустимое» значение (таблица ниже).


\noindent Таблица 3.2 – Классы эквивалентности и граничные значения\par\nopagebreak
\begin{longtable}{|>{\raggedright\arraybackslash}p{0.245\textwidth}|>{\raggedright\arraybackslash}p{0.245\textwidth}|>{\raggedright\arraybackslash}p{0.245\textwidth}|>{\raggedright\arraybackslash}p{0.245\textwidth}|}
\hline
Поле & Допустимые & Недопустимые & Границы \\ \hline
\endhead
name & Буквы 2–60 & Цифры, 1 символ, более 60 & 1 / 2 / 60 / 61 \\ \hline
email & Корректный адрес & Без @, без домена, более 120 & 120 / 121 \\ \hline
team & 0–500 & Отрицательное, больше 500, не число & -1 / 0 / 500 / 501 \\ \hline
notation & Из списка & Любое другое значение & — \\ \hline
comment & 0–500 & Более 500 & 500 / 501 \\ \hline
consent & true & false, отсутствует & — \\ \hline
\end{longtable}
\vspace{\baselineskip}

\reportsection{3.3}{Тест-кейсы и результаты}

Разработано 22 тест-кейсов: позитивные, негативные и граничные. Идентификаторы в таблице сокращены: полный вид TC-WL-NNN. В столбце «Данные» указаны только поля, отличающиеся от базового корректного набора (name=Анна Иванова, email=anna@example.com, notation=idef0, team=5, comment пустой, consent=true). Общие предусловия: проект запущен, форма открыта, поля пусты. Постусловие: данные не сохраняются, состояние системы не меняется. Шаги общие: 1) подать набор данных на вход проверки формы; 2) получить список неверных полей. Ожидаемый результат указан в таблице.


\noindent Таблица 3.3 – Тест-кейсы и результаты выполнения\par\nopagebreak
\begin{longtable}{|>{\raggedright\arraybackslash}p{0.163\textwidth}|>{\raggedright\arraybackslash}p{0.163\textwidth}|>{\raggedright\arraybackslash}p{0.163\textwidth}|>{\raggedright\arraybackslash}p{0.163\textwidth}|>{\raggedright\arraybackslash}p{0.163\textwidth}|>{\raggedright\arraybackslash}p{0.163\textwidth}|}
\hline
ID & Тип, приоритет & Заголовок & Данные & Ожидаемый & Фактический \\ \hline
\endhead
WL-001 & позитивный, High & Все поля корректны & все поля корректны & ошибок нет & ошибок нет (Passed) \\ \hline
WL-002 & граничный, High & Имя из 1 символа (ниже границы 2) & name="A" & name & name (Passed) \\ \hline
WL-003 & граничный, Medium & Имя из 2 символов (нижняя граница) & name="Ян" & ошибок нет & ошибок нет (Passed) \\ \hline
WL-004 & граничный, Medium & Имя из 60 символов (верхняя граница) & name=aaaaa… (60) & ошибок нет & ошибок нет (Passed) \\ \hline
WL-005 & граничный, Medium & Имя из 61 символа (выше границы) & name=aaaaa… (61) & name & name (Passed) \\ \hline
WL-006 & негативный, High & Имя с цифрой & name="Anna1" & name & name (Passed) \\ \hline
WL-007 & позитивный, Low & Имя с апострофом и дефисом & name="O'Neil-Smith" & ошибок нет & ошибок нет (Passed) \\ \hline
WL-008 & негативный, High & Email без символа @ & email=anna.… (16) & email & email (Passed) \\ \hline
WL-009 & негативный, Medium & Email без домена верхнего уровня & email="anna@example" & email & email (Passed) \\ \hline
WL-010 & граничный, Medium & Email длиной 120 символов & email=aaaaa… (120) & ошибок нет & ошибок нет (Passed) \\ \hline
WL-011 & граничный, Medium & Email длиной 121 символ & email=aaaaa… (121) & email & email (Passed) \\ \hline
WL-012 & граничный, High & Размер команды 0 (нижняя граница) & team=0 & ошибок нет & ошибок нет (Passed) \\ \hline
WL-013 & граничный, High & Размер команды 500 (верхняя граница) & team=500 & ошибок нет & ошибок нет (Passed) \\ \hline
WL-014 & граничный, High & Размер команды 501 (выше границы) & team=501 & team & team (Passed) \\ \hline
WL-015 & граничный, Medium & Размер команды -1 (ниже границы) & team=-1 & team & team (Passed) \\ \hline
WL-016 & негативный, Medium & Размер команды — не число & team="abc" & team & team (Passed) \\ \hline
WL-017 & позитивный, Low & Размер команды строкой "12" & team="12" & ошибок нет & ошибок нет (Passed) \\ \hline
WL-018 & негативный, High & Нотация вне списка & notation="xyz" & notation & notation (Passed) \\ \hline
WL-019 & граничный, Low & Комментарий 500 символов & comment=xxxxx… (500) & ошибок нет & ошибок нет (Passed) \\ \hline
WL-020 & граничный, Low & Комментарий 501 символ & comment=xxxxx… (501) & comment & comment (Passed) \\ \hline
WL-021 & негативный, High & Согласие на обработку не дано & consent=false & consent & consent (Passed) \\ \hline
WL-022 & негативный, High & Пустая форма & name=undefined; email=undefined; notation=undefined; team=undefined; comment=undefined; consent=undefined & name, email, notation, team, consent & name, email, notation, team, consent (Passed) \\ \hline
\end{longtable}
\vspace{\baselineskip}

\reportchapter{4}{АНАЛИЗ РЕЗУЛЬТАТОВ}

Выполнено 22 тест-кейсов, пройдено 22, не пройдено 0. Дефектов не выявлено: проверка полей соответствует требованиям на всех границах. Попутно исправлена ошибка самих тестов на этапе отладки (неверный формат тестовых данных), к продукту она отношения не имеет.


\reportchapter{5}{ВЫВОД}

Освоены оформление тест-кейсов по шаблону, эквивалентное разбиение и анализ граничных значений. Тест-кейсы автоматизированы (команда npm run tc:lab2) и повторяются без участия человека.


\end{document}
